\documentclass[11pt,a4paper]{article}
\usepackage{fontspec,amsmath,amssymb,amsthm,microtype}
\setmainfont{lmroman10-regular.otf}[BoldFont=lmroman10-bold.otf,ItalicFont=lmroman10-italic.otf,BoldItalicFont=lmroman10-bolditalic.otf]
\usepackage[margin=25mm]{geometry}
\usepackage[unicode,hidelinks]{hyperref}
\hypersetup{pdfauthor={LangevinLean Project}}
\setlength{\emergencystretch}{3em}
\setlength{\parskip}{0.4em}
\setlength{\parindent}{0pt}
\allowdisplaybreaks[2]
\raggedbottom
\title{\texorpdfstring{Integrating entropy dissipation with almost-everywhere derivatives}{Integrating entropy dissipation with almost-everywhere derivatives}}
\author{LangevinLean Project}
\date{}
\begin{document}\maketitle
Draft mathematical exposition; these notes may contain mistakes.


Suppose the entropy is locally absolutely continuous on positive-time compact intervals and satisfies an affine differential inequality almost everywhere. Multiplication by the exponential integrating factor and the fundamental theorem of calculus give the corresponding affine comparison on each such interval.



At the lower endpoint use an explicit upper bound \(H(s)\le M(s)\), where \(M\) is continuous at zero. Sending \(s\downarrow0\) gives the comparison with \(M(0)\), without requiring continuity or differentiability of \(H\) at zero. Applying this argument to the Fisher dissipation and gradient-moment estimates, then using the stated step restriction, proves the one-step bound with forcing \(6h^2dL^2\).


\section{\texorpdfstring{Integrate the affine differential inequality}{Integrate the affine differential inequality}}


Suppose \(H\) is absolutely continuous on \([s,t]\) and \(H^{\prime}\le-aH+b\) almost everywhere, with \(a>0\). Set \(F(u)=e^{au}(H(u)-b/a)\). Products of absolutely continuous functions on a compact interval are absolutely continuous. The product rule holds almost everywhere and gives \(F^{\prime}(u)=e^{au}(a(H(u)-b/a)+H^{\prime}(u))\le0\).



The fundamental theorem of calculus for absolutely continuous functions gives \(F(t)-F(s)=\int_s^tF^{\prime}\le0\). Multiply by \(e^{-at}>0\), cancel the exponentials, and rearrange:

\[H(t)\le e^{-a(t-s)}H(s)+\frac ba(1-e^{-a(t-s)}).\]

The proof removes no exceptional set by assuming a derivative there; both differentiability and the inequality are used only almost everywhere.



Local absolute continuity gives an integrable derivative and the fundamental theorem on each positive compact interval. Multiplying by the smooth integrating factor preserves absolute continuity, so integrating its almost-everywhere derivative bound is legitimate.


\section{\texorpdfstring{Send the positive initial time to zero}{Send the positive initial time to zero}}


Suppose \(H\) is absolutely continuous on every \([s,T]\), \(0<s<T\), and \(H(s)\le M(s)\) for all sufficiently small positive \(s\). If \(M\) is continuous at zero, the preceding estimate gives

\[H(T)\le e^{-a(T-s)}M(s)+\frac ba(1-e^{-a(T-s)}).\]

The right side tends to \(e^{-aT}M(0)+(b/a)(1-e^{-aT})\). Order is preserved under limits, giving the desired initial-endpoint bound. Single endpoints have zero Lebesgue measure, so an inequality on the open interval suffices. Continuity or differentiability of \(H\) at zero is not needed.



The upper initial majorant converges to the finite initial entropy. Applying it at the lower positive endpoint and sending that endpoint to zero therefore bounds the limiting comparison term from above, which is all the scalar inequality requires.


\section{\texorpdfstring{Absorb the ULA forcing with its exact constants}{Absorb the ULA forcing with its exact constants}}


Assume \(2\alpha H\le I\), \(H^{\prime}\le-3I/4+t^2L^2G_0+2tdL^2\), and \(G_0\le4L^2M(0)/\alpha+2dL\). Let \(0<h\le\alpha/(4L^2)\), \(0<\alpha\le L\), and \(M(0)\ge0\). Then \(hL\le1/4\), \(\alpha h\le1/4\),

\[h^2L^2(4L^2/\alpha)\le\alpha/4,\qquad h^2L^2(2dL)\le hdL^2/2.\]

The derivative is therefore at most \(-3\alpha H/2+B\), where \(B=\alpha M(0)/4+(5/2)hdL^2\). The affine comparison and \(1-e^{-ah}\le ah\) imply \(H(h)\le e^{-3\alpha h/2}M(0)+Bh\). Finally \(e^{-x}\ge e^{-3x/2}+x/4\) for \(0\le x\le1/4\) absorbs the entropy part of \(Bh\). Bounding \(5/2\) by 6 gives

\[H(h)\le e^{-\alpha h}M(0)+6h^2dL^2.\]

The probability-law convergence theorem supplies each of these analytic hypotheses before applying this scalar estimate.



The time-dependent forcing is first bounded uniformly using \(t\le h\). The initial-entropy contribution is absorbed by the exponential gap, while the dimensional forcing integrates to a multiple of \(h^2dL^2\), giving the stated one-step coefficient.


\end{document}
